\subsection{Governing Thermoelastic Equations}

The mathematical formulation of thermoelastic wave propagation is completed by combining the mechanical balance law, the strain--displacement relation, the thermoelastic constitutive relation, and the heat-conduction equation with thermoelastic coupling. These equations define how displacement, strain, stress, and temperature evolve as a connected thermal--mechanical system.

The mechanical motion of the medium is governed by the balance of linear momentum. Neglecting body forces, this balance is the equation of motion already given in Eq.~\eqref{eq:eom}. This equation states that acceleration of material particles is caused by spatial gradients of stress. In the context of wave propagation, a local disturbance produces displacement, displacement produces strain, strain produces stress, and stress gradients drive the motion of the medium.

The strain field is obtained from the displacement field through the infinitesimal strain--displacement relation already introduced in Eq.~\eqref{eq:strain-disp}.
This relation follows from the small-strain assumption and defines the local deformation of the medium. The trace of the strain tensor, \(\varepsilon_{kk}\), represents volumetric strain and is especially important in the thermoelastic coupling term.

The stress field is connected with strain and temperature perturbation through the linear isotropic thermoelastic constitutive relation in Eq.~\eqref{eq:thermoelastic-constitutive}. The first two terms represent the elastic stress caused by mechanical deformation, while the final term represents the thermal stress contribution caused by temperature perturbation. Therefore, temperature change enters the mechanical problem directly through the stress tensor.

The thermal field is described by the coupled heat-conduction equation in Eq.~\eqref{eq:coupled-heat}. The left-hand side describes heat conduction and transient heat storage, while the right-hand side represents coupling with the rate of volumetric deformation. This term expresses the fact that deformation may influence the thermal field in a coupled thermoelastic process.

Together, these equations form the mathematical core of the thermoelastic wave propagation problem. The equation of motion describes mechanical acceleration, the strain--displacement relation connects motion with deformation, the constitutive relation connects deformation and temperature with stress, and the heat equation describes the evolution of temperature. The system is therefore coupled because temperature affects stress, stress affects motion, and volumetric deformation can affect temperature.